% !TEX program = pdflatex
\documentclass[12pt,a4paper]{article}

% -------------------------------------------------
% Packages
% -------------------------------------------------
\usepackage[a4paper,margin=1in]{geometry}
\usepackage[T1]{fontenc}
\usepackage[utf8]{inputenc}
\usepackage[expansion=false]{microtype}
\usepackage{setspace}
\usepackage{parskip}
\usepackage{titlesec}
\usepackage{xcolor}
\usepackage{listings}
\usepackage{fancyhdr}
\usepackage{lastpage}
\usepackage{enumitem}
\usepackage{booktabs}
\usepackage{array}
\usepackage{hyperref}
\usepackage{graphicx}

% -------------------------------------------------
% General formatting
% -------------------------------------------------
\onehalfspacing
\setlength{\parindent}{0pt}

\hypersetup{
    colorlinks=true,
    linkcolor=black,
    urlcolor=blue,
    pdftitle={Blockchain Laboratory - Experiment 2},
    pdfauthor={}
}

\pagestyle{fancy}
\fancyhf{}
\lhead{Blockchain Laboratory}
\rhead{Experiment No. 2}
\cfoot{Page \thepage}
\setlength{\headheight}{15pt}

\titleformat{\section}
  {\large\bfseries}
  {\thesection.}
  {0.5em}
  {}
\titleformat{\subsection}
  {\normalsize\bfseries}
  {\thesubsection.}
  {0.5em}
  {}

% -------------------------------------------------
% Code formatting
% -------------------------------------------------
\definecolor{codebg}{RGB}{248,248,248}
\definecolor{codecomment}{RGB}{80,120,80}
\definecolor{codekeyword}{RGB}{70,70,150}
\definecolor{codestring}{RGB}{150,70,70}

\lstdefinestyle{pythonstyle}{
    language=Python,
    backgroundcolor=\color{codebg},
    basicstyle=\ttfamily\small,
    keywordstyle=\color{codekeyword}\bfseries,
    commentstyle=\color{codecomment},
    stringstyle=\color{codestring},
    numbers=left,
    numberstyle=\tiny\color{gray},
    numbersep=8pt,
    frame=single,
    framerule=0.4pt,
    breaklines=true,
    breakatwhitespace=false,
    showstringspaces=false,
    tabsize=4,
    columns=fullflexible,
    keepspaces=true,
    literate={→}{{$\rightarrow$}}1
}

\lstdefinestyle{outputstyle}{
    backgroundcolor=\color{codebg},
    basicstyle=\ttfamily\small,
    frame=single,
    framerule=0.4pt,
    breaklines=true,
    columns=fullflexible,
    keepspaces=true
}

% -------------------------------------------------
% Document
% -------------------------------------------------
\begin{document}

% -------------------------------------------------
% Title page
% -------------------------------------------------
\begin{titlepage}
    \centering
    \vspace*{1.5cm}

    {\Large\bfseries BLOCKCHAIN LABORATORY\par}
    \vspace{1.2cm}

    {\large\bfseries Experiment No. 2\par}
    \vspace{0.8cm}

    {\LARGE\bfseries Create a Blockchain using Python\par}

    \vspace{1cm}

    \begin{tabular}{rl}
    \textbf{Name:} & Nikhil Sunchu \\[0.6cm]
    \textbf{Roll No.:} & 31 \\[0.6cm]
    \textbf{Class/section} & D20A \\[0.6cm]
    \textbf{Date}& \hrulefill \\[0.6cm]
    \end{tabular}

    \vfill
    {\large Blockchain Laboratory\par}
\end{titlepage}

% -------------------------------------------------
\section{Aim}
Create a Blockchain using Python.

% -------------------------------------------------
\section{Theory}
Blockchain is a distributed and immutable ledger that records data in a sequence of linked blocks. Each block contains a timestamp, proof (from Proof-of-Work), and the hash of the previous block, ensuring integrity and security.

In this program:
\begin{enumerate}[label=\arabic*., leftmargin=*]
    \item \textbf{Block Creation} -- A block is generated using \texttt{create\_block()} with a proof and previous hash.
    \item \textbf{Proof-of-Work (PoW)} -- Implemented to mine a block by solving a cryptographic puzzle (finding a hash with leading zeros).
    \item \textbf{Hashing} -- SHA-256 algorithm ensures data immutability.
    \item \textbf{Validation} -- \texttt{is\_chain\_valid()} checks that every block correctly references the previous hash and satisfies PoW conditions.
    \item \textbf{Flask Web API} -- Routes \texttt{/mine\_block}, \texttt{/get\_chain}, and \texttt{/is\_valid} allow mining, fetching the chain, and verifying blockchain integrity through a browser or API client.
\end{enumerate}

This program runs on a local server and simulates mining a blockchain without a network of nodes, making it an ideal introductory model to understand blockchain basics.

% -------------------------------------------------
\section{Tools and Libraries Used}
\begin{itemize}
    \item Python
    \item Python libraries: \texttt{datetime}, \texttt{hashlib}, \texttt{json}
    \item Flask (web API framework)
\end{itemize}

% -------------------------------------------------
\section{Code Overview}

\subsection{1. Importing Required Libraries}
\begin{itemize}
    \item \texttt{datetime} $\rightarrow$ to timestamp each block
    \item \texttt{hashlib} $\rightarrow$ to generate SHA-256 hashes
    \item \texttt{json} $\rightarrow$ to encode blocks as JSON for hashing
    \item \texttt{Flask} $\rightarrow$ to run the blockchain API server
\end{itemize}

\subsection{2. Blockchain Class -- Core Logic}
This contains everything needed to:
\begin{itemize}
    \item create blocks
    \item store the chain
    \item implement Proof of Work
    \item validate chain integrity
\end{itemize}

\subsection{3. Flask Web API}
Provides endpoints:
\begin{itemize}
    \item \texttt{/mine\_block} $\rightarrow$ mine a new block
    \item \texttt{/get\_chain} $\rightarrow$ fetch full blockchain
    \item \texttt{/is\_valid} $\rightarrow$ verify blockchain
\end{itemize}

% -------------------------------------------------
\section{Complete Code}

\begin{lstlisting}[style=pythonstyle]
# ------------------------------------------------------------
# Module 1 - Create a Simple Blockchain (Fully Commented)
# ------------------------------------------------------------
import datetime
import hashlib
import json
from flask import Flask, jsonify, request

# ------------------------------------------------------------
# Part 1 - Building the Blockchain
# ------------------------------------------------------------
class Blockchain:

    def __init__(self):
        # Store all mined blocks and pending transactions
        self.chain = []
        self.transactions = []

        # Create the genesis block
        self.create_block(
            proof=1,
            previous_hash="0",
            transactions=[]
        )

    def create_block(self, proof, previous_hash, transactions):
        # Build a new block with the current timestamp and transaction list
        block = {
            "index": len(self.chain) + 1,
            "timestamp": str(datetime.datetime.now()),
            "transactions": transactions,
            "proof": proof,
            "previous_hash": previous_hash
        }
        self.chain.append(block)
        return block

    def get_previous_block(self):
        # Return the latest mined block
        return self.chain[-1]

    def hash(self, block):
        # Convert block data to a hash for integrity checking
        encoded_block = json.dumps(
            block,
            sort_keys=True
        ).encode()
        return hashlib.sha256(encoded_block).hexdigest()

    def proof_of_work(self, previous_proof):
        # Find a number whose hash starts with 000 (mining challenge)
        new_proof = 1
        while True:
            hash_operation = hashlib.sha256(
                str(new_proof ** 2 - previous_proof ** 2).encode()
            ).hexdigest()
            if hash_operation[:3] == "000":
                return new_proof
            new_proof += 1

    def is_chain_valid(self, chain):
        # Verify all hashes and proof-of-work values in the chain
        previous_block = chain[0]
        block_index = 1
        while block_index < len(chain):
            block = chain[block_index]
            if block["previous_hash"] != self.hash(previous_block):
                return False

            previous_proof = previous_block["proof"]
            proof = block["proof"]
            hash_operation = hashlib.sha256(
                str(proof ** 2 - previous_proof ** 2).encode()
            ).hexdigest()
            if not hash_operation.startswith("000"):
                return False

            previous_block = block
            block_index += 1
        return True

    def create_transaction(self, sender, receiver, amount):
        # Store a new pending transaction before the next block is mined
        transaction = {
            "sender": sender,
            "receiver": receiver,
            "amount": amount
        }
        self.transactions.append(transaction)
        return len(self.chain) + 1

    def mine_block(self):
        # Mine a new block only when at least one transaction is pending
        if not self.transactions:
            return None

        previous_block = self.get_previous_block()
        previous_proof = previous_block["proof"]
        transactions = self.transactions
        self.transactions = []

        proof = self.proof_of_work(previous_proof)
        previous_hash = self.hash(previous_block)
        block = self.create_block(
            proof=proof,
            previous_hash=previous_hash,
            transactions=transactions
        )
        return block


# ------------------------------------------------------------
# Part 2 - Mining our Blockchain (Flask API)
# ------------------------------------------------------------
app = Flask(__name__)

# Create a blockchain instance for the API server
blockchain = Blockchain()

# ------------------------------------------------------------
# API Endpoints
# ------------------------------------------------------------
@app.route("/get_chain", methods=["GET"])
def get_chain():
    # Return the complete blockchain and its current length
    response = {
        "chain": blockchain.chain,
        "length": len(blockchain.chain)
    }
    return jsonify(response), 200


@app.route("/add_transaction", methods=["POST"])
def add_transaction():
    # Accept a JSON payload and queue a new transaction for mining
    data = request.get_json()
    sender = data["sender"]
    receiver = data["receiver"]
    amount = data["amount"]
    block_index = blockchain.create_transaction(sender, receiver, amount)
    response = {
        "message": "Transaction added successfully",
        "block_index": block_index,
        "transaction": {
            "sender": sender,
            "receiver": receiver,
            "amount": amount
        }
    }
    return jsonify(response), 201


@app.route("/mine_block", methods=["GET"])
def mine_block():
    # Mine the pending transactions into a new block
    block = blockchain.mine_block()
    if block is None:
        return jsonify({
            "message": "No transactions available. Block not mined."
        }), 400
    response = {
        "message": "Block mined successfully.",
        "index": block["index"],
        "timestamp": block["timestamp"],
        "transactions": block["transactions"],
        "proof": block["proof"],
        "previous_hash": block["previous_hash"]
    }
    return jsonify(response), 200


@app.route("/is_valid", methods=["GET"])
def is_valid():
    # Check whether the blockchain has been tampered with
    valid = blockchain.is_chain_valid(blockchain.chain)
    if valid:
        response = {"message": "Blockchain is valid."}
    else:
        response = {"message": "Blockchain is not valid."}
    return jsonify(response), 200


# ------------------------------------------------------------
# Run the Flask app
# ------------------------------------------------------------
if __name__ == "__main__":
    app.run(host="127.0.0.1", port=5000, debug=True)
\end{lstlisting}

% -------------------------------------------------
\section{Output}

\subsection{Fig. 1 -- Get the Initial Chain with Genesis Block}
\begin{center}
\includegraphics[width=0.85\textwidth]{images/fig1.png}
\end{center}

\subsection{Fig. 2 -- Add a Transaction}
\begin{center}
\includegraphics[width=0.85\textwidth]{images/fig2.png}
\end{center}

\subsection{Fig. 3 -- Mine the New Block with the Newly Added Transaction}
\begin{center}
\includegraphics[width=0.85\textwidth]{images/fig3.png}
\end{center}

\subsection{Fig. 4 -- Get the Updated Chain with Mined Block}
\begin{center}
\includegraphics[width=0.85\textwidth]{images/fig4.png}
\end{center}

\subsection{Fig. 5 -- Add Two New Transactions}
\begin{center}
\includegraphics[width=0.85\textwidth]{images/fig5.png}
\end{center}

\subsection{Fig. 6 -- Mine a New Block with the Two Newly Added Transactions}
\begin{center}
\includegraphics[width=0.85\textwidth]{images/fig6.png}
\end{center}

\subsection{Fig. 7 -- Get the Updated Chain with Three Blocks}
\begin{center}
\includegraphics[width=0.85\textwidth]{images/fig7.png}
\end{center}

\subsection{Fig. 8 -- Validate the Chain}
\begin{center}
\includegraphics[width=0.85\textwidth]{images/fig8.png}
\end{center}

% -------------------------------------------------
\section{Conclusion}
We successfully created a basic blockchain using Python and Flask that supports block mining, chain retrieval, and validity checks. Proof-of-Work ensures security by making block creation computationally intensive, while cryptographic hashing guarantees immutability. The implementation demonstrates core blockchain principles -- decentralization, immutability, and consensus -- in a simplified environment, providing a strong foundation for building more advanced blockchain systems.

\end{document}
